% Generated draft: requires mathematical and source review.
\documentclass[11pt]{article}
\usepackage[margin=25mm]{geometry}
\usepackage{fontspec}
\setmainfont{Latin Modern Roman}
\usepackage{amsmath,amssymb,mathtools,unicode-math}
\setmathfont{Latin Modern Math}
\usepackage{enumitem,needspace,xurl,hyperref}
\setlength{\parindent}{0pt}
\setlength{\parskip}{5pt}
\setlength{\emergencystretch}{3em}
\newcommand{\N}{\mathbb{N}}
\newcommand{\Z}{\mathbb{Z}}
\newcommand{\Q}{\mathbb{Q}}
\newcommand{\R}{\mathbb{R}}
\newcommand{\C}{\mathbb{C}}
\newcommand{\F}{\mathbb{F}}
\newcommand{\A}{\mathbb{A}}
\newcommand{\PP}{\mathbb{P}}
\newcommand{\E}{\mathbb{E}}
\newcommand{\eps}{\varepsilon}
\newcommand{\dd}{\,\mathrm{d}}
\begin{document}
\section*{UnsolvedMath ID: 700084}
\textbf{Source problem code:} catalogue entry
\subsection{Definitions and mathematical statement}
Let \(G\) be a group, with no finiteness assumption. A group is locally nilpotent if each of its finitely generated subgroups is nilpotent. Write \(R_{\mathrm{ln}}(G)\) for the locally nilpotent radical of \(G\), namely its largest normal locally nilpotent subgroup.

The supplied record uses Plotkin's terms ``nil-element,'' ``nil-automorphism,'' and ``locally stable automorphism'' without defining them. To retain these unresolved source-specific meanings, write
\[
\begin{aligned}
\mathcal N(G)&=\{g\in G:g\text{ is a nil-element in the source's sense}\},\\
\mathcal A_{\mathrm{nil}}(G)&=\{\alpha\in\operatorname{Aut}(G):\alpha\text{ is a nil-automorphism in the source's sense}\},\\
\mathcal A_{\mathrm{ls}}(G)&=\{\alpha\in\operatorname{Aut}(G):\alpha\text{ is locally stable in the source's sense}\}.
\end{aligned}
\]
These are placeholders for the source terminology, not independent definitions of those three properties.

According to the supplied statement, \(G\) is an \(NR\)-group precisely when
\[
\mathcal N(G)=R_{\mathrm{ln}}(G).
\]
The source also asserts that this condition is equivalent to
\[
\operatorname{Inn}(G)\subseteq\mathcal A_{\mathrm{ls}}(G),
\]
where \(\operatorname{Inn}(G)\) is the group of inner automorphisms of \(G\).

\textbf{Question (V. G. Vilyatser, 1967).} Does there exist a group \(G\) and an automorphism \(\alpha\in\operatorname{Aut}(G)\) such that
\[
\mathcal N(G)=R_{\mathrm{ln}}(G)
\quad\text{and}\quad
\alpha\in\mathcal A_{\mathrm{nil}}(G)\setminus\mathcal A_{\mathrm{ls}}(G)?
\]
A self-contained mathematical formulation requires the missing definitions of the three source-specific properties. The displayed question preserves the logical form of the supplied problem without selecting unconfirmed interpretations.
\subsection{Short English statement}
Vilyatser asks whether an \(NR\)-group can admit a nil-automorphism that is not locally stable. The source characterizes \(NR\)-groups by equality of their nil-elements with their locally nilpotent radical, and equivalently by local stability of every inner automorphism. The meanings of nil-element, nil-automorphism, and local stability require clarification from the source.
\subsection{Sources}
\begin{itemize}
\item Supplied local UnsolvedMath registry record: permanent ID \(700084\), exact source problem code \texttt{catalogue entry}, title ``Kourovka Notebook Problem 2.5.'' Upstream revision: \texttt{d371ee917ac9647fbb06cd6fc529deec6f3dda82}.
\item Source attribution supplied by the registry: \emph{The Kourovka Notebook (Unsolved Problems in Group Theory)}, 2nd issue (1967), problem 2.5; proposed by V. G. Vilyatser. The definition of an \(NR\)-group is attributed to Plotkin.
\item Edition identified in the supplied background: 21st issue edition (2026), \texttt{arXiv:1401.0300v48}. The collection description names E. I. Khukhro and V. D. Mazurov as editors. This bibliographic attribution is supplied source data and has not been independently verified.
\item The registry labels the problem ``open'' and records an update timestamp of \texttt{2026-10-16T14:00:00Z}. This is supplied registry metadata, not an independently verified literature assessment or a stated literature-check date.
\end{itemize}
\end{document}
